\section*{Chapter \ref{ch:g12:risk-and-insurance} --- Risk and Insurance}
% figdata: figdata/grade-12/02-risk-and-insurance/out/ (death_rates, pool_spread); Kestria's orchards invented.

\begin{solution}{exo:g12:risk-and-insurance:1}
$0.2 \times 100 = 20$.
\end{solution}

\begin{solution}{exo:g12:risk-and-insurance:2}
$0.002 \times 150\,000 = 300$ a year.
\end{solution}

\begin{solution}{exo:g12:risk-and-insurance:3}
France 9.1 and 9.9; Japan 11.2 and 11.1; Sweden 8.6 and 9.5; United States 8.7 and 10.3. The American rate rose most, by about 1.6 per 1000; Japan's fell slightly.
\end{solution}

\begin{solution}{exo:g12:risk-and-insurance:4}
\omterm{def:g12:risk-and-insurance:expected-utility}{Expected utility} $(0 + 10)/2 = 5$, so the \omterm{def:g12:risk-and-insurance:certainty-equivalent}{certainty equivalent} is $5^2 = 25$.
\end{solution}

\begin{solution}{exo:g12:risk-and-insurance:5}
$0.03 \times 10\,000 = 300$.
\end{solution}

\begin{solution}{exo:g12:risk-and-insurance:6}
Expected value $(64 + 144)/2 = 104$; \omterm{def:g12:risk-and-insurance:expected-utility}{expected utility} $(8 + 12)/2 = 10$, \omterm{def:g12:risk-and-insurance:certainty-equivalent}{certainty equivalent} 100; \omterm{def:g12:risk-and-insurance:certainty-equivalent}{risk premium} 4.
\end{solution}

\begin{solution}{exo:g12:risk-and-insurance:7}
\omterm{def:g12:risk-and-insurance:expected-utility}{Expected utility} $0.75\sqrt{400} + 0.25\sqrt{100} = 15 + 2.5 = 17.5$; \omterm{def:g12:risk-and-insurance:certainty-equivalent}{certainty equivalent} $17.5^2 = 306.25$. Expected \omterm{def:g11:inequality:wealth}{wealth} $400 - 75 = 325$; \omterm{def:g12:risk-and-insurance:certainty-equivalent}{risk premium} 18.75.
\end{solution}

\begin{solution}{exo:g12:risk-and-insurance:8}
Fair premium $0.25 \times 300 = 75$. The most she pays is $400 - 306.25 = 93.75$.
\end{solution}

\begin{solution}{exo:g12:risk-and-insurance:9}
At 30\,\%: $75 \times 1.3 = 97.5 > 93.75$, she does not buy. At 20\,\%: $90 < 93.75$, she buys.
\end{solution}

\begin{solution}{exo:g12:risk-and-insurance:10}
One member: $1000\sqrt{0.05 \times 0.95} \approx 217.9$. Pool of 100: $217.9/10 \approx 21.8$. Pool of 10\,000: $217.9/100 \approx 2.18$.
\end{solution}

\begin{solution}{exo:g12:risk-and-insurance:11}
Variance $0.25 \times 400 + 0.25 \times 400 = 200$; standard deviation $\sqrt{200} \approx 14.1$ points.
\end{solution}

\begin{solution}{exo:g12:risk-and-insurance:12}
Buyers of home insurance mostly choose low \omterm{def:g11:asymmetric-information:deductible}{deductibles} at costs well above the expected saving. A low \omterm{def:g11:asymmetric-information:deductible}{deductible} insures a small extra loss, over which a concave utility is
nearly linear, so the model predicts near \omterm{def:g12:risk-and-insurance:risk-aversion}{risk neutrality} there; to explain the choices it needs a \omterm{def:g12:risk-and-insurance:risk-aversion}{risk aversion} so high that the same person would refuse almost any
gamble with a large stake, which is implausible. Other explanations, such as overweighting small probabilities or disliking losses, lie outside the model.
\end{solution}

\begin{solution}{exo:g12:risk-and-insurance:13}
$218^2 (0.0001 + 0.1 \times 0.9999) \approx 47\,524 \times 0.10009 \approx 4757$, a standard deviation of about 69, against 2.18 if the losses were independent. As $n$ grows
the variance tends to $\rho\sigma^2$, a standard deviation of $218\sqrt{0.1} \approx 69$: the common part of the risk cannot be pooled away.
\end{solution}

\begin{solution}{exo:g12:risk-and-insurance:14}
The average fair premium is $0.125 \times 64 = 8$. A low risk's \omterm{def:g12:risk-and-insurance:expected-utility}{expected utility} uninsured is $0.95 \times 10 + 0.05 \times 6 = 9.8$, a \omterm{def:g12:risk-and-insurance:certainty-equivalent}{certainty equivalent} of 96.04, so she
pays at most 3.96: she does not buy. If only high risks remain, the fair premium is $0.2 \times 64 = 12.8$; their \omterm{def:g12:risk-and-insurance:certainty-equivalent}{certainty equivalent} is $9.2^2 = 84.64$, so they pay up to
15.36 and buy. Compulsory membership at 8 keeps everyone insured; the low risks pay more than their cover is worth to them, and the high risks gain.
\end{solution}

\begin{solution}{exo:g12:risk-and-insurance:15}
To a risk-neutral grower, the fan lowers the expected loss by $0.05 \times 64 = 3.2$ a year. Fully insured at a fixed premium, he gains nothing from it: his loss is covered
either way. Insurance weakens the incentive to prevent losses unless the premium rewards \omterm{def:g12:risk-and-insurance:prevention}{prevention}: \omterm{def:g11:asymmetric-information:moral-hazard}{moral hazard}.
\end{solution}

\begin{solution}{pb:g12:risk-and-insurance:1}
\textbf{1.} $100 - 0.1 \times 64 = 93.6$.

\textbf{2.} $0.9 \times 10 + 0.1 \times 6 = 9.6$.

\textbf{3.} $9.6^2 = 92.16$.

\textbf{4.} $93.6 - 92.16 = 1.44$.

\textbf{5.} Fair premium $6.4$; at most $100 - 92.16 = 7.84$.

\textbf{6.} $400 \times 6.4 = 2560$ in both.

\textbf{7.} One grower: $64\sqrt{0.1 \times 0.9} = 19.2$. Total: $19.2\sqrt{400} = 384$.

\textbf{8.} Two standard deviations: $2 \times 384 = 768$.

\textbf{9.} $768/400 = 1.92$.

\textbf{10.} For one year, $6.4 + 1.92 = 8.32$ exceeds 7.84: no. But the reserve is kept from year to year; built over two years, the charge is 0.96 and the total $7.36 < 7.84$:
then yes.

\textbf{11.} $0.1 \times 25\,600 = 2560$: unchanged.

\textbf{12.} $400 \times 64 = 25\,600$; the shortfall is $25\,600 - 2560 = 23\,040$.

\textbf{13.} A frost year comes with probability 0.1, more than 2.3\,\%, so the mutual must cover it: 23\,040, thirty times the reserve of Question 8.

\textbf{14.} Every member's loss is the same in each year, so the average loss per member is 64 or 0: adding members adds no independent risk to average out.

\textbf{15.} By passing the valley's risk to an insurer that pools it with risks elsewhere, which are not struck by the same frost; or by a public fund that spreads the
cost across a whole country or across years.

\textbf{16.} It saves $0.05 \times 64 = 3.2$ in expected losses for a cost of 2: yes.

\textbf{17.} An insured grower's losses are paid either way, so the fan saves him nothing: fewer buy it, and the mutual's claims rise. It could require fans, or leave part
of each loss with the grower.

\textbf{18.} A common premium priced on the average risk is above the slope growers' fair premium, possibly above the most they will pay; if they leave, the premium must
rise to the valley-floor risk.

\textbf{19.} Growers who refuse cover at a premium below 7.84, the most a grower with the model's utility would pay; or growers who keep a large share of their losses when
cheap cover is offered.

\textbf{20.} A grower pays at most 7.84 for full insurance (the fair 6.4 plus a \omterm{def:g12:risk-and-insurance:certainty-equivalent}{risk premium} of 1.44). To pay all claims in at least 97.7\,\% of years, the mutual needs a
reserve of 768 beyond its premiums with independent frosts, and 23\,040 with a common frost.
\end{solution}
